\chapter{1977年湖北卷}

一、计算

\begin{problem}
\begin{enumerate}[label=（\arabic*）]
\item \( \sqrt{(-3)^2} \)；\( \sqrt[3]{(-27)^3} \).

\item \( \log_2 2+\log_5 1-\log_{\frac{1}{3}}\left( \frac{1}{27} \right) \).

\item \( \left( \sqrt{\frac{2}{3}} \right)^3 \cdot \sqrt[3]{\left( \frac{3}{2} \right)^2} \cdot \left( \frac{2}{3} \right)^{-\frac{5}{6}} \).

\item 解方程：\( x-\sqrt{8-x^2}=0 \).

\item 当 \( m\le 1 \) 时，计算 \( \sqrt{(m-1)^2} \).

\item 求函数 \( y=\frac{\sqrt{2-x}}{x+4} \) 的定义域.

\item 分解因式：\( x^4+x^2y^2+y^4 \).

\item 证明：\( \tan 22^\circ30'=\sqrt{2}-1 \)（不查表）.
\end{enumerate}

\end{problem}

\begin{solution}
（1）\( \sqrt{(-3)^2}=\sqrt{3^2}=3 \)；

\( \sqrt[3]{(-27)^3}=-27 \).

（2）
\[
\log_2 2+\log_5 1-\log_{\frac{1}{3}}\left( \frac{1}{27} \right)
=1+0-\log_{\frac{1}{3}}\left( \left( \frac{1}{3} \right)^3 \right)
=1-3\log_{\frac{1}{3}}\frac{1}{3}
=1-3=-2.
\]

（3）
\[
\left( \sqrt{\frac{2}{3}} \right)^3 \cdot \sqrt[3]{\left( \frac{3}{2} \right)^2} \cdot \left( \frac{2}{3} \right)^{-\frac{5}{6}}
=\sqrt[6]{\left( \frac{2}{3} \right)^9} \cdot \sqrt[6]{\left( \frac{3}{2} \right)^4} \cdot \sqrt[6]{\left( \frac{3}{2} \right)^5}
\]
\[
=\sqrt[6]{\left( \frac{2}{3} \right)^9 \cdot \left( \frac{3}{2} \right)^9}=1.
\]

（4）移项，得 \( x=\sqrt{8-x^2} \).

两边平方，得 \( x^2=8-x^2 \).

所以 \( x^2=4 \)，\( x=\pm 2 \).

经检验，知 \( x=-2 \) 不适合原方程，是增根舍去，故原方程的解为 \( x=2 \).

（5）因为 \( m\le 1 \)，所以 \( m-1\le 0 \)，故
\[
\sqrt{(m-1)^2}=1-m.
\]

（6）要使函数 \( y=\frac{\sqrt{2-x}}{x+4} \) 有意义，自变量 \( x \) 必须满足下面的不等式：
\[
\left\{
\begin{aligned}
2-x&\ge 0,\\
x+4&\ne 0.
\end{aligned}
\right.
\]
解之，得 \( -4<x\le 2 \) 和 \( x<-4 \).

因此，函数的定义域为 \( -4<x\le 2 \) 和 \( x<-4 \) 的实数（也可表示为 \( x\le 2 \) 但 \( x\ne -4 \) 的实数）.

（7）在有理式范围内可作如下因式分解：
\[
x^4+x^2y^2+y^4=\left( x^2+y^2 \right)^2-x^2y^2
=\left( x^2+y^2-xy \right)\left( x^2+y^2+xy \right).
\]

（8）
\[
\tan 22^\circ30'=\tan\frac{45^\circ}{2}=\frac{1-\cos45^\circ}{\sin45^\circ}
=\frac{1-\frac{\sqrt{2}}{2}}{\frac{\sqrt{2}}{2}}
=\frac{2-\sqrt{2}}{\sqrt{2}}
=\frac{2\left( \sqrt{2}-1 \right)}{2}
=\sqrt{2}-1.
\]
\end{solution}

\begin{problem}
围一矩形场地，一边利用房屋的墙，其他三边用 \( L \) 米的篱笆围成，问怎样围才能使面积最大？最大面积是多少？

\bitmapfigure[max width=4.2cm,max totalheight=3.0cm]{img/1977/hubei/q2_fig1.png}

\end{problem}

\begin{solution}
如图，设围成的矩形场地的长为 \( x \) 米，那么宽就应是 \( \frac{L-x}{2} \) 米.

这时矩形面积
\[
S=x\left( \frac{L-x}{2} \right)=-\frac{1}{2}x^2+\frac{L}{2}x
\]
\[
=-\frac{1}{2}\left( x-\frac{L}{2} \right)^2+\frac{L^2}{8}.
\]

所以当 \( x=\frac{L}{2} \)（米）时，\( S \) 有最大值 \( \frac{L^2}{8} \)（米\(^2\)）. 这时的宽 \( \frac{L-x}{2}=\frac{L}{4} \)（米）.

答：当与墙平行的一边长为 \( \frac{L}{2} \) 米，与墙垂直的两边分别长为 \( \frac{L}{4} \) 米时，围成的矩形面积最大. 最大面积为 \( \frac{L^2}{8} \) 米\(^2\).
\end{solution}

\begin{problem}
已知直线 \( L_1 \) 方程为 \( x+y-1=0 \) 和点 \( A(-3,0) \).

\begin{enumerate}[label=（\arabic*）]
\item 求与直线 \( L_1 \) 垂直且通过 \( A \) 的直线方程；

\item 求点 \( A \) 到直线 \( L_1 \) 的距离. 
\end{enumerate}

\end{problem}

\begin{solution}
（1）直线 \( L_1 \) 的斜率为 \( -1 \)，因此和 \( L_1 \) 垂直的直线的斜率为 \( 1 \). 根据点斜式，所求的直线方程为
\[
\frac{y-0}{x+3}=1,
\]
或者写成 \( x-y+3=0 \).

（2）点 \( A \) 到直线 \( L_1 \) 的距离
\[
d=\frac{|-3+0-1|}{\sqrt{1+1}}=\frac{4}{\sqrt{2}}=2\sqrt{2}.
\]
\end{solution}

\begin{problem}
解方程组：\(\left\{
\begin{aligned}
\log_2 x-\log_4 y&=0,\\
y^2-13x^2+36&=0
\end{aligned}
\right.\).

\end{problem}

\begin{solution}
原方程组可写成：
\[
\left\{
\begin{aligned}
\log_2 x-\frac{1}{2}\log_2 y&=0,\\
y^2-13x^2+36&=0.
\end{aligned}
\right.
\]

即
\[
\left\{
\begin{aligned}
x&=\sqrt{y}\quad \text{（1）},\\
y^2-13x^2+36&=0\quad \text{（2）}.
\end{aligned}
\right.
\]

把（1）代入（2），得
\[
y^2-13y+36=0.
\]
\[
(y-9)(y-4)=0.
\]

所以 \( y_1=9 \)，\( y_2=4 \).

把 \( y_1=9 \)，\( y_2=4 \) 分别代入（1）式，求得 \( x_1=3 \)，\( x_2=2 \). 经检验，都是原方程的根. 所以原方程的解为
\[
\left\{
\begin{aligned}
x_1&=3,\\
y_1&=9
\end{aligned}
\right.
\quad \text{和} \quad
\left\{
\begin{aligned}
x_2&=2,\\
y_2&=4.
\end{aligned}
\right.
\]
\end{solution}

\begin{problem}
作函数 \( y=\sqrt{3}\cos 2x+\sin 2x \) 的图象（要求写出作图步骤，并画出草图）.

\end{problem}

\begin{solution}
\[
y=\sqrt{3}\cos 2x+\sin 2x
=2\left( \frac{\sqrt{3}}{2}\cos 2x+\frac{1}{2}\sin 2x \right)
\]
\[
=2\left( \sin\frac{\pi}{3}\cos 2x+\cos\frac{\pi}{3}\sin 2x \right)
=2\sin\left( 2x+\frac{\pi}{3} \right).
\]

所以这个函数的周期为 \( \pi \)，振幅为 \( 2 \).

\bitmapfigure[max width=6.4cm,max totalheight=8.8cm]{img/1977/hubei/q5_solution_fig1.png}
\end{solution}

\referenceproblemsection
\setcounter{probnum}{0}

\begin{problem}
求下列函数的导数：

\( y=\sin x \)；\( y=\e^{3x} \)；\( y=\ln x \)；\( f(t)=A\sin\left( \omega t+\varphi \right) \).

\end{problem}

\begin{solution}
\( y=\sin x \)，则 \( y'=\cos x \).

\( y=\e^{3x} \)，则 \( y'=3\e^{3x} \).

\( y=\ln x \)，则 \( y'=\frac{1}{x} \).

\( f(t)=A\sin\left( \omega t+\varphi \right) \)，则 \( f'(t)=\omega A\cos\left( \omega t+\varphi \right) \).
\end{solution}

\begin{problem}
设函数 \( f(x)=x^3-12x+10 \).

\begin{enumerate}[label=（\arabic*）]
\item 求 \( f(x) \) 的极值点；

\item 求 \( f(1) \)，\( f(2) \)，\( f(-3) \).
\end{enumerate}

\end{problem}

\begin{solution}
（1）在极值点，必定 \( f'(x)=0 \).

即 \( f'(x)=3x^2-12=0 \).

解得 \( x=\pm 2 \). 这时 \( f(2)=-6 \)，\( f(-2)=26 \). 所以函数 \( f(x)=x^3-12x+10 \) 的极值点为 \( (2,-6) \) 和 \( (-2,26) \).

（2）\( f(1)=1-12+10=-1 \).

\( f(2)=2^3-12\times 2+10=-6 \).

\( f(-3)=(-3)^3-12\times(-3)+10=19 \).
\end{solution}

\begin{problem}
解方程：\( y''+3y'-4y=0 \)，其中 \( y|_{x=0}=0 \)，\( y'|_{x=0}=-10 \).

\end{problem}

\begin{solution}
\( y''+3y'-4y=0 \) 是线性齐次方程，它的特征方程 \( \lambda^2+3\lambda-4=0 \) 的根是 \( -4 \) 和 \( 1 \). 因此方程的通解是
\[
y=C_1\e^{-4x}+C_2\e^x
\]
（\( C_1 \)，\( C_2 \) 是任意常数）.

代入初值问题，\( y|_{x=0}=0 \) 和 \( y'|_{x=0}=-10 \)，得方程组：
\[
\left\{
\begin{aligned}
C_1+C_2&=0,\\
-4C_1+C_2&=-10.
\end{aligned}
\right.
\]

求得 \( C_1=2 \)，\( C_2=-2 \). 故得初值问题的解是
\[
y=2\e^{-4x}-2\e^x.
\]
\end{solution}
